% TOP_PROBLEM: {"id": 30007382, "title": "A first-order temperature singularity for a stable radial continuum gas", "category_id": 19, "category": {"id": 19, "name": "probability", "display_name": "Probability", "description": "Problems involving probability theory, stochastic processes, and random structures.", "slug": "probability", "order_index": 19, "created_at": "2026-07-31T15:26:25.671Z"}, "status": "open"}
% FIRST_POSTED: 2026-10-06
% ATTEMPT_STATUS: solved
% OPENAI_CLAIM: {"schema_version":1,"source_repo":"https://github.com/openai/math","source_commit":"adc7f1241b42e322a6451854ab7e4b4c146bf78a","release_date":"2026-10-06","match":"full","resolution":"proved","independently_reviewed":false,"second_pass":"agreed","families":[{"family":"228","title":"Continuum phase transitions for radial pair potentials","subject":"Probability and statistical mechanics","manuscripts":[{"title":"A continuum temperature singularity for a radial pair potential","pdf_path":"preprints/A-continuum-temperature-singularity-for-a-radial-pair-potential-September-24-2026/paper.pdf","readme_path":"preprints/A-continuum-temperature-singularity-for-a-radial-pair-potential-September-24-2026/README.md","date":"2026-09-24","theorem_ref":"thm:main (introduction.tex:44)","statement_sources":[{"path":"preprints/A-continuum-temperature-singularity-for-a-radial-pair-potential-September-24-2026/build/introduction.tex","line":44,"label":"thm:main"}]},{"title":"A radial continuum phase transition with algebraic decay","pdf_path":"preprints/A-radial-continuum-phase-transition-with-algebraic-decay-September-24-2026/paper.pdf","readme_path":"preprints/A-radial-continuum-phase-transition-with-algebraic-decay-September-24-2026/README.md","date":"2026-09-24","theorem_ref":"thm:main (sections/introduction.tex:32)","statement_sources":[{"path":"preprints/A-radial-continuum-phase-transition-with-algebraic-decay-September-24-2026/build/sections/introduction.tex","line":32,"label":"thm:main"}]}],"lean":{"doc_path":"lean/docs/228.md","comparators":["lean/ComparatorChallenges/ContinuumTransition.json","lean/ComparatorChallenges/ContinuumTransition.lean","lean/ComparatorChallenges/RadialDensityInterval.json","lean/ComparatorChallenges/RadialDensityInterval.lean","lean/ComparatorChallenges/RadialTransition.json","lean/ComparatorChallenges/RadialTransition.lean"],"declarations":["OAI.ContinuumTemperature.main_theorem","OAI.RadialTransition.main"],"covers_main_theorem":true},"reasoning_trace":null}]}
\documentclass[11pt]{article}
\usepackage[margin=25mm]{geometry}
\usepackage{fontspec}
\setmainfont{Latin Modern Roman}
\usepackage{amsmath,amssymb,mathtools}
\usepackage{unicode-math}
\setmathfont{Latin Modern Math}
\usepackage{enumitem}
\usepackage{xurl,hyperref,bookmark}
\hypersetup{colorlinks=true,urlcolor=blue}
\setcounter{secnumdepth}{0}
\setlength{\parindent}{0pt}
\setlength{\parskip}{5pt}
\setlength{\emergencystretch}{3em}
\begin{document}
\section*{A first-order temperature singularity for a stable radial continuum gas}

\subsection{OpenAI's claim}
The linked manuscripts present the following claims by OpenAI: A stable three-dimensional radial pair potential with a divergent core and integrable tail is claimed to exhibit one finite-temperature canonical derivative jump throughout a positive density interval; A bounded continuous stable radial potential with decay $O(r^{-3-1/32})$ is claimed to give the same common-temperature first-order singularity over an open density interval. Together they give a claimed affirmative answer in the matching scope.

\subsection{Mathematical statement of the claimed result}
\paragraph{Result 1.} There exists a radial Borel pair potential $\Phi(x)=\phi(|x|)$ on $\mathbb R^3$ with $\phi(0)=+\infty$, finite locally bounded values on $(0,\infty)$, and
\[
 \phi(r)\to+\infty\ (r\downarrow0),\quad \phi(r)\le-a\ (r\in[r_1,r_2]),\quad
 \int_R^\infty r^2|\phi(r)|\,dr<\infty,\quad \phi(r)=o(r^{-3}),
\]
for suitable $a>0$, $0<r_1<r_2$ and finite $R$. Its interaction satisfies $U_N=\sum_{i<j}\Phi(x_i-x_j)\ge-BN$ on every finite configuration. For free cubes $\Lambda_L=[-L/2,L/2]^3$, let
\[
 Z_{L,N}(\beta)=\frac1{N!}\int_{\Lambda_L^N}e^{-\beta U_N}\,dx_1\cdots dx_N,\qquad
 f_\Phi(\beta,\rho)=-\beta^{-1}\lim_{L\to\infty}L^{-3}\log Z_{L,\lfloor\rho L^3\rfloor}(\beta).
\]
This limit exists and is finite for every $\beta,\rho>0$. For some nonempty open density interval $I$ and a single $\beta_c\in(1/2,3/2)$,
\[
 D^-_\beta f_\Phi(\beta_c,\rho)>D^+_\beta f_\Phi(\beta_c,\rho)\quad(\rho\in I),
\]
with finite one-sided derivatives. Position integration uses temperature-independent Lebesgue measure.

\paragraph{Result 2.} A second construction gives a bounded continuous radial pair potential $\phi$ on $\mathbb R^3$ and finite $B,C>0$ satisfying
\[
 \sum_{i<j}\phi(|x_i-x_j|)\ge-BN,\qquad |\phi(r)|\le Cr^{-3-1/32}\quad(r\ge1).
\]
Let $p=\lim_{L\to\infty}L^{-3}M_L$, where $M_L$ is the maximum number of points in $[0,L]^3$ with pairwise separation at least one. For the free-cube partition function
\[
 Z_{L,N}^\phi(\beta)=\frac1{N!}\int_{([0,L]^3)^N}
 e^{-\beta\sum_{i<j}\phi(|x_i-x_j|)}\,dx_1\cdots dx_N,
\]
the limit $f(\beta,\rho)=-\beta^{-1}\lim_LL^{-3}\log Z_{L,\lfloor\rho L^3\rfloor}^\phi(\beta)$ is finite for every $\beta>0$ and every $\rho$ in an open interval $I$ centred at $5p/3$. At a common $\beta_c\in[7/8,9/8]$,
\[
 f'_-(\beta_c,\rho)>f'_+(\beta_c,\rho)\quad(\rho\in I),
\]
with both derivatives finite.

\subsection{Outline of the argument}
Four-particle cluster bounds give stability and canonical thermodynamic limits. Successive Kac-type attractions preserve a triple of supporting slopes; perturbation estimates keep the slope separation while the range grows, yielding one potential and one common singular temperature.

Packing bounds control cell occupations and ensure stability of the bounded core. Long-range attraction is compared quantitatively with an ideal density transform, and pressure supports are preserved through the range iteration; the chosen scales produce the explicit algebraic tail.

\subsection{Scope relative to this problem}
The first-order singularity is constructed for its admissible potentially unbounded-core class with only o(r\textasciicircum{}(-3)) tail and beta\_c\(\in\)(1/2,3/2). It does not establish the target bounded-continuous r\textasciicircum{}(-3\(-\)1/32) potential with narrower critical interval. A bounded continuous stable radial potential with decay $O(r^{-3-1/32})$ is claimed to give the same common-temperature first-order singularity over an open density interval. The question retains these same parameters and conventions. Within these assumptions, this is a claimed affirmative answer to the stated question.

\subsection{Formal verification}
For family 228, the documented constructions include both the divergent-core singularity and the bounded continuous algebraically decaying potential, with one common inverse temperature over their respective density intervals. The listed Comparator files are \texttt{ContinuumTransition.lean}, \texttt{RadialDensityInterval.lean}, \texttt{RadialTransition.lean}. Named declarations include \texttt{OAI.ContinuumTemperature.main\_theorem}, \texttt{OAI.RadialTransition.main}. This site has not executed the Lean code. The scope descriptions do not certify the remaining manuscript arguments.

\subsection{Status caveat}
OpenAI attributes these manuscripts to work by an internal model. Its release cautions that arguments without formalization can contain mistakes. This catalogue has not independently checked the proofs; the labels report the mathematical claims attributed to OpenAI.

\subsection{Sources}
\begin{itemize}
\item OpenAI, \emph{A continuum temperature singularity for a radial pair potential}, 2026-09-24; citation key \texttt{OAI:A-continuum-temperature-singularity-for-a-radial-pair-potential-September-24-2026}. \href{https://github.com/openai/math/blob/adc7f1241b42e322a6451854ab7e4b4c146bf78a/preprints/A-continuum-temperature-singularity-for-a-radial-pair-potential-September-24-2026/paper.pdf}{PDF (GitHub)}; \href{https://github.com/openai/math/raw/adc7f1241b42e322a6451854ab7e4b4c146bf78a/preprints/A-continuum-temperature-singularity-for-a-radial-pair-potential-September-24-2026/paper.pdf}{PDF (direct)}; \href{https://github.com/openai/math/blob/adc7f1241b42e322a6451854ab7e4b4c146bf78a/preprints/A-continuum-temperature-singularity-for-a-radial-pair-potential-September-24-2026/README.md}{Manuscript page and citation}; \href{https://github.com/openai/math/blob/main/preprints/A-continuum-temperature-singularity-for-a-radial-pair-potential-September-24-2026/paper.pdf}{Latest version}; Statement source: \href{https://github.com/openai/math/blob/adc7f1241b42e322a6451854ab7e4b4c146bf78a/preprints/A-continuum-temperature-singularity-for-a-radial-pair-potential-September-24-2026/build/introduction.tex\#L44}{thm:main}.
\item OpenAI, \emph{A radial continuum phase transition with algebraic decay}, 2026-09-24; citation key \texttt{OAI:A-radial-continuum-phase-transition-with-algebraic-decay-September-24-2026}. \href{https://github.com/openai/math/blob/adc7f1241b42e322a6451854ab7e4b4c146bf78a/preprints/A-radial-continuum-phase-transition-with-algebraic-decay-September-24-2026/paper.pdf}{PDF (GitHub)}; \href{https://github.com/openai/math/raw/adc7f1241b42e322a6451854ab7e4b4c146bf78a/preprints/A-radial-continuum-phase-transition-with-algebraic-decay-September-24-2026/paper.pdf}{PDF (direct)}; \href{https://github.com/openai/math/blob/adc7f1241b42e322a6451854ab7e4b4c146bf78a/preprints/A-radial-continuum-phase-transition-with-algebraic-decay-September-24-2026/README.md}{Manuscript page and citation}; \href{https://github.com/openai/math/blob/main/preprints/A-radial-continuum-phase-transition-with-algebraic-decay-September-24-2026/paper.pdf}{Latest version}; Statement source: \href{https://github.com/openai/math/blob/adc7f1241b42e322a6451854ab7e4b4c146bf78a/preprints/A-radial-continuum-phase-transition-with-algebraic-decay-September-24-2026/build/sections/introduction.tex\#L32}{thm:main}.
\item Family 228: \href{https://github.com/openai/math/blob/adc7f1241b42e322a6451854ab7e4b4c146bf78a/lean/docs/228.md}{Lean scope}.
\item \href{https://github.com/openai/math/blob/adc7f1241b42e322a6451854ab7e4b4c146bf78a/lean/ComparatorChallenges/ContinuumTransition.json}{Comparator ContinuumTransition.json}.
\item \href{https://github.com/openai/math/blob/adc7f1241b42e322a6451854ab7e4b4c146bf78a/lean/ComparatorChallenges/ContinuumTransition.lean}{Comparator ContinuumTransition.lean}.
\item \href{https://github.com/openai/math/blob/adc7f1241b42e322a6451854ab7e4b4c146bf78a/lean/ComparatorChallenges/RadialDensityInterval.json}{Comparator RadialDensityInterval.json}.
\item \href{https://github.com/openai/math/blob/adc7f1241b42e322a6451854ab7e4b4c146bf78a/lean/ComparatorChallenges/RadialDensityInterval.lean}{Comparator RadialDensityInterval.lean}.
\item \href{https://github.com/openai/math/blob/adc7f1241b42e322a6451854ab7e4b4c146bf78a/lean/ComparatorChallenges/RadialTransition.json}{Comparator RadialTransition.json}.
\item \href{https://github.com/openai/math/blob/adc7f1241b42e322a6451854ab7e4b4c146bf78a/lean/ComparatorChallenges/RadialTransition.lean}{Comparator RadialTransition.lean}.
\item \href{https://github.com/openai/math/blob/adc7f1241b42e322a6451854ab7e4b4c146bf78a/CONTENTS.md}{OpenAI catalogue entry}.
\end{itemize}
\end{document}
